\problemname{Binary Palindromes}
\noindent

Mikael is investigating interesting positive integers, and has now stumbled upon \emph{binary palindromes}. To formally describe this, we define the following functions:

\[
\text{bin}(n) = \text{the binary representation of the integer } n \text{ without leading zeros},
\]

\[
\text{reverse}(s) = \text{the string } s \text{ written backwards}.
\]

An integer $n$ is therefore a \emph{binary palindromes} if
\[
\text{bin}(n) = \text{reverse}(\text{bin}(n)).
\]

Note that a palindrome is a string that reads the same forwards and backwards.

Mikael likes the binary palindromes that contain exactly $k$ ones in their binary representation, and now wants to know how many such numbers exist in the interval $[l, r]$.

Given three integers $l$, $r$, and $k$, determine how many integers $x$ satisfy:
\begin{enumerate}
    \item The integer $x$ is between $l$ and $r$, i.e.\ $l \leq x \leq r$,
    \item The integer $x$ is a number Mikael likes (i.e.\ a binary palindrome),
    \item The binary representation of $x$ contains exactly $k$ ones.
\end{enumerate}

\section*{Input}
The input consists of a single line containing three integers $l, r, k$ $(1 \leq l \leq r \leq 10^{18}, 1 \leq k < 60)$.

\section*{Output}
Output a single integer, the number of integers in the interval $[l, r]$ that are binary palindromes and contain exactly $k$ ones in their binary representation.

\section*{Scoring}
Your solution will be tested on a set of test groups, each worth a number of points. 
Each test group contains a set of test cases. 
To get the points for a test group you need to solve all test cases in the test group.

\noindent
\begin{tabular}{| l | l | p{12cm} |}
  \hline
  \textbf{Group} & \textbf{Points} & \textbf{Constraints} \\ \hline
  $1$    & $20$       & $l,r \leq 10^9$ \\ \hline
  $2$    & $20$       & $k \leq 4$ \\ \hline
  $3$    & $60$       & No additional constraints \\ \hline
\end{tabular}

\section*{Explanation of Sample 1}

We seek numbers in the interval $[7, 10]$ that are binary palindromes and have exactly two ones in their binary representation.

There are only 4 numbers in the interval:

\begin{align*}
7 &= 111_2 \quad (\text{palindrome, three ones}) \\ 
8 &= 1000_2 \quad (\text{not a palindrome}) \\
9 &= 1001_2 \quad (\text{palindrome, two ones}) \\
10 &= 1010_2 \quad (\text{not a palindrome})
\end{align*}

The only number satisfying both conditions is $9$, so the answer is $1$.

\section*{Explanation of Sample 2}

We seek numbers in the interval $[1, 6]$ that are binary palindromes and have exactly two ones in their binary representation.

There are 6 numbers in the interval:

\begin{align*}
1 &= 1_2 \quad (\text{palindrome, a single one}) \\
2 &= 10_2 \quad (\text{not a palindrome}) \\
3 &= 11_2 \quad (\text{palindrome, two ones}) \\
4 &= 100_2 \quad (\text{not a palindrome}) \\
5 &= 101_2 \quad (\text{palindrome, two ones}) \\
6 &= 110_2 \quad (\text{not a palindrome})
\end{align*}

The numbers satisfying both conditions are $3$ and $5$, so the answer is $2$.
